\chapter{Compartment Syndrome}
\index{Compartment Syndrome}
\label{sec:Compartment Syndrome}

\section{Overview}
Compartment syndrome is elevated pressure within a closed fascial compartment, compromising tissue perfusion and leading to reduced blood flow and tissue death of muscle and nerves. Acute compartment syndrome is a surgical emergency most commonly caused by fractures (tibial fractures most common), crush injuries, reperfusion injury, or tight casts. This genetic disorder results from specific gene mutations or chromosomal abnormalities. It may be present at birth or manifest later in life depending on the underlying mechanism. Genetic disorders result from abnormalities in an individual's DNA, ranging from single-gene mutations to chromosomal abnormalities. These conditions may be inherited from parents or arise from new mutations. Pattern of inheritance depends on whether the mutation is dominant, recessive, or X-linked. Genetic counseling helps families understand risks and make informed decisions. This condition affects a significant number of people worldwide and can range from mild to severe in presentation. Early recognition and appropriate management are essential for optimal outcomes. A thorough understanding of the underlying mechanisms guides treatment decisions. Patient education and self-management play important roles in long-term care.
\section{Causes}
The underlying mechanisms involve complex interactions between genetic predisposition and environmental factors. Disruption of normal physiological processes leads to the characteristic features of the condition. Multiple contributing factors may act synergistically to trigger or worsen the condition.
\section{Risk Factors}
Family history of the genetic condition Consanguineous parents (increased risk of recessive disorders) Advanced parental age (new mutations) Carrier status in parents Ethnic background (specific founder mutations) Previous child with a genetic disorder

\begin{itemize}[noitemsep]
  \item Genetic predisposition and family history
  \item Advancing age
  \item Lifestyle factors (diet, exercise, smoking, alcohol)
  \item Environmental exposures
  \item Underlying medical conditions
  \item Medication use
\end{itemize}
\section{Signs and Symptoms}
\begin{itemize}[noitemsep]
  \item You may experience pain out of proportion to the injury, worsened by passive stretching of the affected muscles, with tense, swollen compartments
  \item The five P's (pulselessness, tingling or numbness, paralysis, pallor, poikilothermia) are late signs.
  \item Pain or discomfort in the affected area
  \item Functional impairment affecting daily activities
  \item Changes in appearance or function of affected tissues
  \item Systemic symptoms may include fatigue, fever, or weight changes
  \item Symptoms may develop gradually or appear suddenly
\end{itemize}
\section{Progression and Stages}
The condition may progress through different stages of severity over time. Early stages often involve mild or subtle symptoms that may be overlooked. Moderate stages involve more noticeable symptoms affecting daily function. Advanced stages may involve significant impairment and complications.
\section{Complications}
\begin{itemize}[noitemsep]
  \item Disease progression leading to worsening symptoms
  \item Secondary infections or injuries
  \item Side effects of treatment
  \item Reduced quality of life
  \item Emotional and psychological impact
\end{itemize}
\section{Diagnosis}
Doctors diagnose this based on clinical examination and compartment pressure measurement (pressures >30 mmHg or within 30 mmHg of diastolic pressure indicate need for fasciotomy).

\begin{itemize}[noitemsep]
  \item Comprehensive medical history and physical examination
  \item Laboratory tests to assess organ function
  \item Imaging studies to visualize affected structures
  \item Specialized testing depending on the affected system
  \item Consultation with relevant specialists
\end{itemize}
\section{Prevention}
\begin{itemize}[noitemsep]
  \item Maintain a healthy lifestyle with balanced diet and exercise
  \item Avoid known risk factors
  \item Attend regular medical check-ups for early detection
  \item Follow recommended screening guidelines
\end{itemize}
\section{Treatment Options}
\begin{itemize}[noitemsep]
  \item Treatment typically involves emergent fasciotomy to decompress all affected compartments
  \item Delayed treatment leads to irreversible muscle tissue death, contracture (Volkmann's ischemic contracture), and potentially limb loss.
  \item Medications to manage symptoms and address underlying mechanisms
  \item Lifestyle modifications and self-care measures
  \item Physical therapy and rehabilitation
  \item Surgical intervention when indicated
  \item Regular monitoring and follow-up care
\end{itemize}
\section{Living with It}
\begin{itemize}[noitemsep]
  \item Follow your treatment plan as prescribed
  \item Maintain regular follow-up with your healthcare provider
  \item Make necessary lifestyle adjustments
  \item Seek support from family, friends, or support groups
  \item Stay informed about your condition
\end{itemize}
\section{Outlook}
Most people recover well with early fasciotomy but poor if Treatment typically involves delayed.
